% Reviewer response: the G2 degradation attributed to "heavier communication
% traffic and coexistence effects".
%
% Explanation, not results. The argument follows from the transport being
% broadcast and from the structure of the coupling term, and is stated without
% measured values in the body. The measurements live in Figure fig:completeness
% and its caption; the full tables are in docs/PLATFORM.md 5.5b.
%
% Figure panels are generated from runs/n30-campaign, 60 cycles per arm at both
% publication rates, n ~ 600 node-runs per cell:
%     python3 tools/capture_analysis.py runs/n30-campaign --plot results/capture
% then copy results/capture/capture_completeness_{8hz,5hz}.png into figures/.
%
% The argument in one line: airtime does not scale with the degree, but the
% RECEIVER's task does. It must acquire d values per coordination interval from
% d uncoordinated transmitters, and the probability of doing so falls
% geometrically in d at fixed airtime and fixed per-link quality.

\begin{enumerate}
\item The degraded performance under $\mathcal{G}_2$ is attributed to heavier
communication traffic and coexistence effects. Please provide quantitative
evidence for this attribution.

\textbf{Response:} The audit prompted by this question changed our explanation.
The traffic argument was inherited from the previous version of the platform and
does not carry over to the current one. We replace it with a mechanism that is
specific to broadcast reception and that does scale with the degree.

\textbf{1. Transmitted airtime does not depend on the degree.} On the revised
platform both media are broadcast, so an agent transmits its state once per
publication period irrespective of how many neighbours read it: airtime per node
is $O(1)$ in the degree rather than $O(|\mathcal{N}_i|)$. Per-node transmission
counts are recorded with each run, and per-link reception is reconstructed at the
receiver from a sequence number carried in the payload, so the claim is checked
rather than assumed.

The previous version of the platform retrieved each neighbour's state with a
separate HTTP request per neighbour per tick. Its traffic genuinely scaled with
the degree, and the original attribution was sound for that implementation. Under
broadcast it is not.

\textbf{2. Received airtime does not depend on the degree either.} A BLE scanner
listens to the channel, not to a neighbour. It occupies the antenna for
$W_\mathrm{scan}/T_\mathrm{scan}$ of the time and reports every advertisement it
decodes, whether or not the sender is one of its neighbours, so its receive
occupancy is a configuration constant: identical for a degree-2 and a degree-4
node in the same manifest. The degree changes which fraction of what was already
received is useful, not how long the receiver listens. The controller counters
distinguish the two, counting every reception event and then classifying it per
neighbour after decoding.

Consequently the coexistence exposure on the Raspberry Pi is also independent of
the degree. That exposure is $W_\mathrm{scan}/T_\mathrm{scan}$ for reception plus
the advertising duty for transmission, and neither term contains
$|\mathcal{N}_i|$.

\textbf{3. What the degree does change: the receiver's acquisition problem.} This
is the mechanism we now advance, and it is a property of reception rather than of
traffic.

\emph{On BLE there is no scheduler.} Every agent advertises on its own interval
and its own phase; the interval is deliberately randomised within a range to
prevent phases from locking. Nothing coordinates the order in which neighbours
transmit, and nothing tells a receiver when to expect a given neighbour's next
advertisement, so reception there is an uncoordinated capture problem rather
than a scheduled read. On Wi-Fi it is not: the same agents are served by an
access point which buffers broadcast frames and releases them on its own beacon
boundaries, and that is a scheduler. The two media therefore differ in exactly
the variable this section is about while running the same law, over the same
degrees, in the same fleet, at the same instant. Figure~\ref{fig:completeness}
is that comparison.

\emph{Per-link capture is degree-independent; the neighbourhood is not.} Whether
a particular neighbour's advertisement is captured depends on the scan duty
cycle, on whether the receiver was itself transmitting at that instant, on
collisions with other advertisers within range, and on the phase relationship
between the two devices. None of these quantities contains the receiver's degree.
Write $\rho$ for the resulting per-link capture probability, which is thus the
same for a sparse and a dense node in the same deployment.

The coordination input, however, is not a function of one link. It is evaluated
over the whole neighbourhood, so within a given publication interval the receiver
must capture $d = |\mathcal{N}_i|$ separate advertisements, from $d$ transmitters
that are not coordinated with each other or with it. Treating the captures as
independent, the probability that the coupling term is computed from a fully
refreshed neighbourhood is
\begin{equation}\label{eq:complete}
    P_\mathrm{complete}(d) = \rho^{\,d},
\end{equation}
which decays geometrically in $d$ even though $\rho$, the airtime and the duty
cycle are all unchanged.

\emph{Both sides of \eqref{eq:complete} are counted, not inferred.} Each agent
logs, at every local update, a flag per neighbour marking that a fresh value
from that neighbour was written into its table. Samples are grouped into
publication windows, and within a window a neighbour counts as acquired if at
least one fresh value from it arrived anywhere in that window. The per-link
capture $\rho$ is then the fraction of windows in which a given neighbour was
acquired, averaged over links, and $P_\mathrm{complete}$ is the fraction of
windows in which all $d$ were. The second is a joint count over the
neighbourhood rather than a product of the first, so $\rho^{\,d}$ is a
prediction tested against an independent measurement rather than a restatement
of it. Grouping by window rather than by sample is deliberate: the per-sample
flag under-counts when two advertisements land inside one sampling period,
whereas whether at least one arrived within a window does not depend on that.

\emph{The error also grows, not only its likelihood.} The expected number of
stale terms in the sum is $d\,(1-\rho)$, linear in the degree. So a denser graph
is more often working from incomplete information \emph{and} the incompleteness
is larger when it occurs. The two effects act in the same direction.

\emph{Redundancy separates per-link from per-event quality, and confirms
independence a second way.} Publication and advertising run on independent
timers, so each published value is radiated
$k = T_\mathrm{pub}/T_\mathrm{adv}$ times before the controller overwrites it.
If advertising events are lost independently with per-event capture probability
$p$, then $\rho = 1-(1-p)^{k}$. Taking $p$ from the configuration with $k = 1$
predicts $\rho$ for the configuration with $k = 1.6$ to within two points on the
BLE agents, at an unchanged radio configuration. Losses are therefore memoryless
from one advertising event to the next, which is the signature of a
channel-contention process rather than of a scheduling artefact, and which is
what makes redundancy a usable lever: raising $k$ raises $\rho$, and
\eqref{eq:complete} amplifies that gain geometrically in the degree.

\begin{figure}[t]
  \centering
  \includegraphics[width=\linewidth]{figures/capture_completeness_8hz.png}\\[2pt]
  \includegraphics[width=\linewidth]{figures/capture_completeness_5hz.png}
  \caption{Neighbourhood completeness against degree at both publication rates,
  with agents grouped by the medium on which they receive. $\mathcal{G}_1$
  supplies the degree-1 point and $\mathcal{G}_2$ the degree-4 point.
  (a)~Per-link capture $\rho$ is flat in the degree, which is the premise of the
  argument. (b)~Measured completeness (solid) against the independent prediction
  $\rho^{d}$ (dotted), annotated with the ratio between them. The BLE agents,
  whose transmitters nothing sequences, lie on their own $\rho^{d}$ curve, so
  \eqref{eq:complete} is what those receivers do rather than an idealisation.
  The Wi-Fi agents, served by an access point that batches broadcast frames, lie
  far above it: correlated arrival makes a complete neighbourhood more frequent
  than independence predicts, and \eqref{eq:complete} is a lower bound there.
  Bridges receive on both media and sit between the two. (c)~Stale neighbours per
  publication window against the prediction $d(1-\rho)$.}
  \label{fig:completeness}
\end{figure}

\textbf{4. Summary of the revised attribution.} Under broadcast, a denser graph
costs nothing in transmitted or received airtime and nothing in coexistence
exposure, provided the access point is configured off the advertising channels
and dual-transceiver modules arbitrate between their radios, so the original
explanation does not apply to this platform. What it does cost is the
reliability of the coordination input: the receiver must acquire more values per
interval, from more uncoordinated transmitters, within the same listening time,
and the probability of assembling a complete neighbourhood falls geometrically
in the degree. The degradation under $\mathcal{G}_2$ is therefore attributed to
reception completeness rather than to channel load.

The single piece of evidence we would point to is the split between the two
media, shown in Figure~\ref{fig:completeness}. They run the same law over the
same degrees in the same fleet, and the one whose transmitters are sequenced by
an access point loses far less to the degree than the one whose transmitters are
sequenced by nothing. Channel load does not distinguish those two cases; the
acquisition argument does.

\end{enumerate}
